\section*{Hoofdstuk \ref{ch:b2:isotope-geochemistry} --- Isotopengeochemie}

\begin{solution}{exo:b2:isotope-geochemistry:1}
Tristan: $(0.51249/0.512638 - 1)\times10^4 = -2.9$. Mangaia: $(0.51285/0.512638 -
1)\times10^4 = +4.1$. MORB: $+8.5$. De bronnen van Mangaia en van MORB, met positieve
$\varepsilon_{\mathrm{Nd}}$, hebben gemiddeld een hogere Sm/Nd gehad dan \omterm{def:b2:planet-formation-meteorites:chondrule}{chondrieten}; de
bron van Tristan een lagere.
\end{solution}

\begin{solution}{exo:b2:isotope-geochemistry:2}
$\delta^{18}\mathrm{O}_{\mathrm{VPDB}} = (\num{2.0672e-3}/\num{2.0052e-3} - 1)\times1000 =
+30.9$ promille op de VSMOW-schaal. Het calciet heeft $^{18}\mathrm{O}/^{16}\mathrm{O} =
\num{2.0672e-3}\times(1 - 0.0020) = \num{2.0631e-3}$, dus $\delta^{18}\mathrm{O} =
(\num{2.0631e-3}/\num{2.0052e-3} - 1)\times1000 = +28.9$ promille (\omterm{def:b2:isotope-geochemistry:standards}{VSMOW}).
\end{solution}

\begin{solution}{exo:b2:isotope-geochemistry:3}
Bij \qty{1.0}{Ga}: $e^{\lambda t} - 1 = 0.006545$, en $0.512638 - 0.1966\times0.006545 =
0.511351$. Bij \qty{4.567}{Ga}: $e^{\lambda t}-1 = 0.030243$, en $0.512638 -
0.1966\times0.030243 = 0.506692$.
\end{solution}

\begin{solution}{exo:b2:isotope-geochemistry:4}
$\delta\mathrm{D} = 8\times(-18.0) + 10 = -134$ promille. Voor het meer is $d = -42 -
8\times(-4.0) = -10$ promille, onder de 10 van de lijn: het water ligt onder de \omterm{prop:b2:isotope-geochemistry:mwl}{meteorische
waterlijn}, zoals verdampte oppervlaktewateren (de lichte isotopen vertrokken eerst).
\end{solution}

\begin{solution}{exo:b2:isotope-geochemistry:5}
$Q = 10^4\times\num{6.524e-12}\times0.1966/0.512638 = \num{2.502e-8}$ per jaar, of
\qty{25.0}{\per\giga\year}. Het graniet heeft $\mathrm{Sm}/\mathrm{Nd} = 6.0/30 = 0.200$,
tegen $0.145/0.457 = 0.317$ voor CI: $f = 0.200/0.317 - 1 = -0.37$. Over
\qty{0.5}{Gyr}: $\Delta\varepsilon = 25.0\times(-0.37)\times0.5 = -4.6$ eenheden.
\end{solution}

\begin{solution}{exo:b2:isotope-geochemistry:6}
$\varepsilon \approx 3.5 + 25.0\times(-0.45)\times2.7 = 3.5 - 30.4 = -26.9$, tegen
$-27.3$: de benadering is 0.4 eenheid te hoog (niet negatief genoeg). Het grootste deel van het
verschil, ongeveer 0.3 eenheid, komt van $e^{\lambda t} > 1$: de moederkern, en dus het
groeitekort ten opzichte van CHUR, was in het verleden overvloediger dan vandaag
(\cref{exo:b2:isotope-geochemistry:10}); de rest komt van het afronden van $f$ en de beginwaarde.
\end{solution}

\begin{solution}{exo:b2:isotope-geochemistry:7}
Met een massafractie $x$ graniet is $\mathrm{Sr} = 320x + 130(1-x)$ en
$R = (320x\times0.7200 + 130(1-x)\times0.7030)/\mathrm{Sr}$. Voor $x = 0.05$: Sr
\qty{139.5}{ppm}, $R = 0.70495$; $x = 0.10$: 149.0, 0.70665; $x = 0.20$: 168.0,
0.70948. Volgens het bewijs van de mengpropositie is $R\,\mathrm{Sr}$ lineair in Sr:
$R = 0.731632 - 3.7221/\mathrm{Sr}$, een rechte lijn met helling $-3.722$ (ppm) en
snijpunt 0.7316 tegen $1/\mathrm{Sr}$. Controle bij $x = 0.05$: $0.731632 - 3.7221/139.5
= 0.70495$.
\end{solution}

\begin{solution}{exo:b2:isotope-geochemistry:8}
Bij \qty{20}{\degreeCelsius}: $\delta^{18}\mathrm{O} = 1000/1.0098 - 1000 = -9.70$
promille, $\delta\mathrm{D} = 1000/1.084 - 1000 = -77.5$ promille, helling $77.5/9.70 = 7.99$.
Bij \qty{0}{\degreeCelsius}: $-11.56$ en $-99.9$ promille, helling 8.64. Beide liggen dicht bij
de helling 8 van de \omterm{prop:b2:isotope-geochemistry:mwl}{meteorische waterlijn}: de evenwichtsfractionering van water
tussen vloeistof en damp alleen brengt de verhouding van de verschuivingen van deuterium en zuurstof-18
al in de buurt van 8.
\end{solution}

\begin{solution}{exo:b2:isotope-geochemistry:9}
De korst bevat $0.0051\times27/1.25 = 0.110$, 11 procent van het neodymium.
Massabalans: $0.110\times(-15) + 0.890\,\varepsilon_m = 0$, dus $\varepsilon_m = 1.65/0.890 =
+1.9$. Met een verarmd reservoir dat een aandeel $x$ bevat op $+8.5$ en de rest op 0:
$0.110\times(-15) + 8.5x = 0$, $x = 1.65/8.5 = 0.19$: ongeveer een vijfde van het
neodymium.
\end{solution}

\begin{solution}{exo:b2:isotope-geochemistry:10}
De groeisnelheid is $\lambda p\,e^{\lambda t}$, met $e^{\lambda t} \approx 1 + \lambda t$;
gemiddeld over het interval van de afsplitsing $T$ tot het heden is de factor
$1 + \lambda T/2$, dus de verwaarloosde termen hebben een relatieve grootte $\lambda T/2$. Over de
ouderdom van de aarde is dat $\num{6.524e-12}\times\num{4.567e9}/2 = 0.015$, 1.5 procent. Voor de
korst van \cref{exo:b2:isotope-geochemistry:6} is $\lambda T/2 = 0.0088$, en de fout is $0.0088\times30.4 = 0.27$
eenheid. De trage verandering van de CHUR-verhouding waardoor $\varepsilon$ wordt gedeeld, voegt
nog een kleinere term toe, van de orde $(p/R)\lambda T \times \varepsilon_0$, hier 0.02 eenheid.
\end{solution}

\begin{solution}{exo:b2:isotope-geochemistry:11}
$K = 13.3/5.0 = 2.66$. Met $s = 0.5$: $s' = 2.66\times0.5/(1 + 1.66\times0.5) =
1.33/1.83 = 0.73$. Het verschil $s' - s$ is het grootst waar
$\mathrm{d}s'/\mathrm{d}s = K/(1+(K-1)s)^2 = 1$, bij $s = (\sqrt K - 1)/(K-1) =
0.631/1.66 = 0.38$, waar $s' = 0.62$: 0.24 verschil.
\end{solution}

\begin{solution}{exo:b2:isotope-geochemistry:12}
Vorming: $2.3$ tot $2.6\times\qty{7.2e9}{km^3}/\qty{4.5e9}{yr} = 3.7$ tot
\qty{4.2}{km^3} per jaar. Vernietiging: $1.3$ tot $1.6\times7.2/4.5 = 2.1$ tot
\qty{2.6}{km^3} per jaar. De huidige snelheden verschillen: de toevoeging (2.7) ligt onder de gemiddelde
vorming van het model, de verwijdering (3.2) boven de gemiddelde vernietiging. In het model groeide de korst
vroeg, toen de vorming het snelst was, en is het heden een tijd van evenwicht of licht
verlies, zoals de balans van de subductiezones zegt.
\end{solution}

\begin{solution}{pb:b2:isotope-geochemistry:1}
\textbf{1.} Mantel: $(0.51307/0.512638 - 1)\times10^4 = +8.4$. Sediment:
$(0.51141/0.512638 - 1)\times10^4 = -24.0$.

\textbf{2.} Kora $+2.8$, Vel $-4.3$, Asha $+4.7$.

\textbf{3.} Met $^{87}\mathrm{Sr}/^{86}\mathrm{Sr}$ horizontaal en $\varepsilon_{\mathrm{Nd}}$ verticaal:
mantel (0.70282, $+8.4$), Kora (0.70410, $+2.8$), Vel (0.70618, $-4.3$), sediment
(0.72000, $-24.0$), Asha (0.70290, $+4.7$). Alleen Vel, met een negatieve
$\varepsilon_{\mathrm{Nd}}$, ligt aan de verrijkte kant van CHUR.

\textbf{4.} Onder de reeks: bij de strontiumverhouding van de \omterm{def:b2:isotope-geochemistry:reservoirs}{verarmde mantel} zou de reeks
een $\varepsilon_{\mathrm{Nd}}$ rond $+8$ hebben; Asha heeft $+4.7$. Het lijkt op de
St.-Helenagroep van het HIMU-type (Mangaia, Tubuai, St.~Helena).

\textbf{5.} Smelten verandert de verhoudingen moeder--dochter, maar niet meetbaar de
isotopenverhoudingen van strontium en neodymium (propositie over het smelten); de lava's zijn
jong, dus hun verhoudingen zijn sinds de eruptie niet gegroeid (vraag 6).

\textbf{6.} $\Delta R = 0.12\,(e^{\lambda\cdot10^6} - 1) = 0.12\times\num{6.52e-6} =
\num{7.8e-7}$, of $\num{7.8e-7}/0.512638\times10^4 = 0.015$ epsiloneenheid: ver onder de
verschillen tussen de eilanden, en onder de precisie van een meting.

\textbf{7.} $\mathrm{Nd} = 40f + 0.75(1-f)$, en
$R = \bigl(40f\times0.51141 + 0.75(1-f)\times0.51307\bigr)/\mathrm{Nd}$.

\textbf{8.} $f = 0.5$ procent: Nd \qty{0.9463}{ppm}, $\varepsilon = +1.6$; 1 procent:
1.1425, $-2.9$; 2 procent: 1.535, $-8.4$; 5 procent: 2.7125, $-15.4$.

\textbf{9.} Sr 10.95, 11.90, 13.80 en \qty{19.50}{ppm}; $^{87}\mathrm{Sr}/^{86}\mathrm{Sr} =
0.70439$, 0.70571, 0.70780 en 0.71163.

\textbf{10.} $R\,\mathrm{Sr} = 10\times0.70282 + f(200\times0.72000 - 10\times0.70282)$ en
$\mathrm{Sr} = 10 + 190f$; na eliminatie van $f$ is $R\,\mathrm{Sr} = 7.0282 +
(\mathrm{Sr}-10)\times136.9718/190$, dus $R = 0.720904 - 0.18084/\mathrm{Sr}$: een rechte
lijn met helling $-0.1808$ (ppm) tegen $1/\mathrm{Sr}$. Controle bij $f = 1$ procent:
$0.720904 - 0.18084/11.90 = 0.70571$.

\textbf{11.} $K = (10/0.75)/(200/40) = 13.33/5.0 = 2.67$. De mantel is relatief
rijker aan strontium, het sediment aan neodymium, dus de Nd-verhouding beweegt sneller naar het
sediment: een beetje sediment verschuift $\varepsilon_{\mathrm{Nd}}$ ver voordat
$^{87}\mathrm{Sr}/^{86}\mathrm{Sr}$ veel beweegt.

\textbf{12.} 53 keer rijker aan Nd ($40/0.75$), 20 keer aan Sr. Bij $f = 1$ procent levert het
sediment $0.40/1.1425 = 35$ procent van het neodymium, wat
$\varepsilon_{\mathrm{Nd}}$ met $0.35\times(8.4 + 24.0) = 11.3$ eenheden verschuift, maar slechts $2.0/11.9 =
17$ procent van het strontium.

\textbf{13.} $0.01\times40/1.1425 = 0.350$: 35 procent.

\textbf{14.} Uit $R\,(40f + 0.75(1-f)) = 40fR_S + 0.75(1-f)R_D$:
$f = 0.75(R - R_D)/\bigl(40(R_S - R) + 0.75(R - R_D)\bigr)$, en hetzelfde met 200
en 10 voor strontium.

\textbf{15.} $f = 0.75\times(-0.00029)/\bigl(40\times(-0.00137) + 0.75\times(-0.00029)\bigr)
= \num{-2.175e-4}/\num{-5.502e-2} = 0.0040$: 0.40 procent.

\textbf{16.} $f = 10\times0.00128/\bigl(200\times0.01590 + 10\times0.00128\bigr) =
0.0128/3.1928 = 0.0040$: 0.40 procent; de twee komen overeen.

\textbf{17.} Neodymium: $0.75\times(-0.00065)/\bigl(40\times(-0.00101) +
0.75\times(-0.00065)\bigr) = \num{-4.875e-4}/\num{-4.089e-2} = 0.0119$; strontium:
$10\times0.00336/\bigl(200\times0.01382 + 10\times0.00336\bigr) = 0.0336/2.7976 = 0.0120$.
Beide geven 1.2 procent.

\textbf{18.} $f = 0.75\times(-0.00019)/\bigl(40\times(-0.00147) + 0.75\times(-0.00019)\bigr)
= 0.0024$, wat Sr \qty{10.46}{ppm} en $^{87}\mathrm{Sr}/^{86}\mathrm{Sr} =
0.70361$ voorspelt, tegen 0.70290 gemeten. Asha ligt niet op de mengkromme mantel--sediment:
zijn bron bevat nog een andere component.

\textbf{19.} Oude gerecyclede oceanische korst, het HIMU-type: zijn Nd-isotopen vragen een Sm/Nd
iets onder die van de \omterm{def:b2:isotope-geochemistry:reservoirs}{verarmde mantel}, zijn Sr-isotopen een lage Rb/Sr, en zijn
lood (hier niet gemeten) een hoge U/Pb, alles over één tot twee miljard jaar.

\textbf{20.} $e^{\lambda_{147}\cdot2\times10^9} - 1 = 0.013133$. Sediment: $0.51141 -
0.110\times0.013133 = 0.509965$; CHUR: $0.512638 - 0.1966\times0.013133 = 0.510056$;
$\varepsilon_{\mathrm{Nd}}(2\,\mathrm{Ga}) = (0.509965/0.510056 - 1)\times10^4 = -1.8$.

\textbf{21.} $f = 0.110/0.1966 - 1 = -0.440$; sinds de subductie veranderde $\varepsilon$ met
$25.0\times(-0.440)\times2.0 = -22.0$, dus $\varepsilon(2\,\mathrm{Ga}) \approx -24.0 + 22.0
= -2.0$, binnen 0.2 eenheid van de exacte waarde.

\textbf{22.} $e^{\lambda_{87}\cdot2\times10^9} - 1 = 0.028338$;
$0.72000 - 0.50\times0.028338 = 0.70583$.

\textbf{23.} Bij de subductie had het sediment $\varepsilon_{\mathrm{Nd}} = -1.8$, dicht bij de
$-4.7$ van een korst die toen pas 0.7 miljard jaar oud was, en ruim boven de
sterk negatieve waarden van oude korst: het past bij de erosie van jonge korst, misschien
gemengd met wat materiaal uit de mantel. Zijn $^{87}\mathrm{Sr}/^{86}\mathrm{Sr}$ van
0.7058 op dat moment is ook die van een jong korstgesteente.

\textbf{24.} De bron van Vel, het meest verrijkte eiland, bevat
\textbf{\qty[text-series-to-math=true]{1.2}{\percent}} gerecycled sediment.
\end{solution}
